\chapter{Metodologia} \label{cap:proposta}

    A metodologia adotada nesta tese estrutura-se de forma progressiva, articulando modelagem matemática, implementação computacional e análise de desempenho com vistas ao refinamento do formalismo proposto. Essa progressão organiza-se em duas etapas: uma primeira fase de verificação de viabilidade, já concluída e com resultados preliminares obtidos, dedicada à construção e validação dos Hamiltonianos propostos; e uma segunda fase, em desenvolvimento, voltada à implementação de algoritmos de otimização quântica variacional sobre os formalismos estabelecidos.

    Inicialmente, o trabalho concentrou-se na construção de Hamiltonianos capazes de representar problemas de otimização combinatória em um contexto de computação quântica de variáveis contínuas. Essa etapa teve como propósito estabelecer uma correspondência entre a estrutura do problema e sua codificação em operadores compatíveis com sistemas quânticos de variáveis contínuas.

    O TSP foi adotado, nesse contexto, como problema base para validação, por constituir uma instância clássica de otimização combinatória com estrutura suficientemente rica para testar o método, mas com complexidade ainda tratável do ponto de vista analítico e computacional. Como ponto de partida, foi considerada uma primeira abordagem baseada em uma correspondência direta entre a modelagem discreta usual e o formalismo de variáveis contínuas. Nesse contexto, cada qubit, originalmente associado a uma variável binária no modelo clássico, foi mapeado para um qumode, de modo a reproduzir o comportamento discreto esperado.

    Para isso, foi construído um Hamiltoniano análogo ao utilizado em formulações do tipo QUBO em qubits. Considera-se um conjunto de $N$ cidades em um grafo, com $d_{ij}$ representando a distância entre as cidades $i$ e $j$, e define-se a variável binária

    $$
    x_{i,k} =
    \begin{cases}
    1, & \text{se a cidade $i$ é visitada na posição $k$},\\
    0, & \text{caso contrário}.
    \end{cases}
    $$

    A função objetivo do TSP pode então ser escrita como:

    $$
        \label{eq:objective_function_hamiltonian}
        H = \sum_{k=1}^N \sum_{i=1}^N \sum_{j=1}^N d_{ij}\, x_{i,k} x_{j,k+1}
    $$

    Sujeito às restrições:

    \begin{enumerate}
        \item Cada posição deve conter exatamente uma cidade:
        $$
            \sum_{i=1}^N x_{i,k} = 1, \qquad \forall k
        $$

        \item Cada cidade deve ser visitada exatamente uma vez:
        $$
            \sum_{k=1}^N x_{i,k} = 1, \qquad \forall i
        $$
    \end{enumerate}

    No contexto de variáveis contínuas, as variáveis discretas $x_{i,k}$ são substituídas por operadores $\hat{x}_{i,k}$ associados aos qumodes. Introduzem-se, adicionalmente, os operadores canônicos de posição $\hat{x}_{i,k}$ e momento $\hat{p}_{i,k}$, satisfazendo as relações de comutação

    $$
    [\hat{x}_{i,k}, \hat{p}_{j,l}] = i\hbar\,\delta_{ij}\delta_{kl}.
    $$
    \noindent adota-se, para simplificação algébrica, a convenção $\hbar=1$. Uma vez que o operador posição possui espectro contínuo, impõe-se um comportamento efetivamente binário por meio da transformação


    $$
    \hat{x}_{i,k} = \frac{1 + \hat{x}'_{i,k}}{2},
    $$

    \noindent de modo que $\hat{x}'_{i,k} \approx \pm 1$ corresponde a $\hat{x}_{i,k} \in \{0,1\}$. Essa restrição é imposta energeticamente através de um potencial de poço duplo:

    $$
    H_{\text{bin}} = \sum_{i,k} \left( (\hat{x}'_{i,k})^2 - 1 \right)^2.
    $$

    Além disso, considera-se um termo cinético dado por:

    $$
    \sum_{i,k} \frac{\hat{p}_{i,k}^2}{2m}.
    $$

    Dessa forma, o Hamiltoniano completo do sistema é definido como:

    \begin{equation}
    \begin{aligned}
    H &= \sum_{k=1}^N \sum_{i=1}^N \sum_{j=1}^N d_{ij}\, \hat{x}_{i,k}\hat{x}_{j,k+1}
        + \alpha \sum_{i,k} \left( (\hat{x}'_{i,k})^2 - 1 \right)^2 \\
        &\quad + \beta \sum_{i=1}^N \left(\sum_{k=1}^N \hat{x}_{i,k} - 1 \right)^2
        + \gamma \sum_{k=1}^N \left(\sum_{i=1}^N \hat{x}_{i,k} - 1 \right)^2
        + \sum_{i,k} \frac{\hat{p}_{i,k}^2}{2m},
    \end{aligned}
    \end{equation}

    onde $\alpha$, $\beta$, $\gamma$ e $m$ são hiperparâmetros do modelo. A solução do problema é obtida pela minimização do valor esperado do Hamiltoniano:

    \begin{equation}
    \min_{\psi \in \mathcal{H}} \bra{\psi} H \ket{\psi}.
    \end{equation}

    Essa formulação inicial teve como principal objetivo verificar a consistência da representação com variáveis contínuas em relação aos modelos tradicionais. Para isso, os resultados obtidos foram comparados tanto com soluções provenientes de uma busca exaustiva clássica quanto com implementações do Hamiltoniano equivalente em qubits, formulação da qual o modelo de qumodes deriva diretamente, por substituição das variáveis binárias por operadores de posição com comportamento binário imposto energeticamente. Em ambos os casos, os estados de energia mínima encontrados no modelo com qumodes coincidiram com os das abordagens de referência, indicando que a modelagem proposta preserva corretamente a estrutura do problema.

    Os testes parciais foram conduzidos para $N \in \{3, 4\}$, onde $N$ denota a quantidade de vértices. Para $N = 3$, o algoritmo de força bruta identificou o circuito hamiltoniano de menor custo como $V_0 \rightarrow V_1 \rightarrow V_2 \rightarrow V_0$, resultado compatível com o obtido pelo algoritmo quântico, que produziu o caminho $V_0 \rightarrow V_2 \rightarrow V_1 \rightarrow V_0$, sendo ambos equivalentes a menos de ordem e com custos mínimos iguais. Analogamente, para $N = 4$, o algoritmo de força bruta retornou o caminho $V_0 \rightarrow V_1 \rightarrow V_2 \rightarrow V_3 \rightarrow V_0$, ao passo que o algoritmo quântico produziu $V_0 \rightarrow V_3 \rightarrow V_2 \rightarrow V_1 \rightarrow V_0$, novamente com custos mínimos idênticos. 

    Cabe ressaltar, contudo, que tais experimentos foram restritos a instâncias de pequena dimensão, em razão do elevado custo computacional associado à simulação de sistemas de variáveis contínuas em hardware clássico. Para um sistema composto por $M$ qubits, a representação explícita do vetor de estados requer memória que cresce como $O(2^M)$; já no caso de $M$ qumodes truncados na base de Fock com corte $d$, esse custo escala como $O(d^M)$. Como a dimensão efetiva do espaço de Hilbert depende diretamente de $d$, a simulação clássica de sistemas contínuos torna-se rapidamente proibitiva à medida que $M$ cresce. Nesse contexto, por instâncias de pequena dimensão entendem-se aquelas em que os valores de $M$ e de $d$ permitem armazenar e manipular numericamente os estados e operadores do sistema de forma computacionalmente viável.

    Como sequência da proposta, foi desenvolvida uma segunda versão do Hamiltoniano, na qual a representação da rota é realizada por meio de variáveis inteiras $n_k \in \{1,2,\dots,N\}$, em que o índice $k$ identifica o vértice $V_k$ e o valor $n_k$ indica a posição em que esse vértice é visitado no percurso. Dessa forma, a configuração do percurso é descrita pelo vetor $(n_1, n_2, \dots, n_N)$, cuja $k$-ésima entrada corresponde ao vértice $V_k$ e cujo valor indica sua ordem de visita na rota. Para exemplificar, considere um estado da forma $\ket{2413}$: o vértice $V_1$ é visitado na segunda posição, $V_2$ na quarta, $V_3$ na primeira e $V_4$ na terceira, resultando na rota $V_3\rightarrow V_1\rightarrow V_4\rightarrow V_2$.

    Nesse contexto, o Hamiltoniano de custo associado ao percurso é dado por:
    $$
    H = \sum_{k=1}^N d_{n_k, n_{k+1}}.
    $$

    Seja $\hat{n}_k$ o operador de número do modo $k$, definido por:
    $$
    \hat{n}_k\ket{a} = a\ket{a},
    $$

    \noindent pode-se escrever o Hamiltoniano correspondente à função objetivo na forma:
    $$
    H = \sum_{k=1}^N \sum_{a=1}^N \sum_{b=1}^N d_{ab}\, \ket{a}\bra{a}_k \otimes \ket{b}\bra{b}_{k+1}.
    $$

    Para impedir a ocorrência de cidades repetidas na rota, introduz-se um termo de penalização definido por:
    $$
    H_{\text{penalty}} = \sum_{k<l} \sum_{a=1}^N \ket{a}\bra{a}_k \otimes \ket{a}\bra{a}_l.
    $$

    Esse termo atribui energia positiva a configurações nas quais uma mesma cidade ocupa múltiplas posições na rota, enquanto estados que correspondem a permutações válidas não são penalizados. Alternativamente, essa penalização pode ser reescrita por meio do operador:
    $$
    \hat{N}_a = \sum_{k=1}^N \ket{a}\bra{a}_k,
    $$

    \noindent resultando na forma equivalente:
    $$
    H_{\text{penalty}} = \sum_{a=1}^N \left(\hat{N}_a - 1 \right)^2.
    $$

    {\color{red} Alterar o parágrafo abaixo, inserindo a tabela de resultados}

    Os resultados preliminares obtidos com essa modelagem foram comparados com os resultados provenientes do mesmo algoritmo de força bruta clássico utilizado na primeira versão do Hamiltoniano, o qual permite enumerar todas as rotas possíveis. Devido ao elevado custo computacional associado à simulação de sistemas com múltiplos qumodes, essa análise foi restrita a instâncias de pequena dimensão. Observou-se que as rotas de menor custo identificadas pelo modelo coincidem com aquelas obtidas pela abordagem clássica, indicando a consistência da formulação proposta.

    Todas as simulações descritas foram executadas em um servidor, denominado Beluga, equipado com processador AMD EPYC 7453 de 56 núcleos e 112 threads, 1 TB de memória RAM e duas GPUs NVIDIA RTX A5500 de 24 GB cada, sob o sistema operacional Ubuntu 24.04.3 LTS. A implementação dos Hamiltonianos e a simulação dos sistemas quânticos foram realizadas com o auxílio da biblioteca QuTiP (\textit{Quantum Toolbox in Python}) \cite{qutip5}, que fornece estruturas para a representação e manipulação de operadores e estados quânticos em espaços de Hilbert de dimensão finita. Os algoritmos utilizados ao longo deste projeto estão apresentados no \autoref{sec:algorithm_append}. 

    Como extensão imediata do formalismo para o caso de múltiplos veículos, propõe-se associar a cada vértice $V_k$ não apenas a variável inteira $n_k \in \{1,2,\dots,N\}$, que indica sua posição de visita, mas também uma variável discreta $v_k \in \{1,2,\dots,M\}$, que identifica o veículo responsável por seu atendimento, sendo $M$ o número total de veículos disponíveis. Desse modo, cada configuração passa a ser descrita pelo estado:
    $$
    \ket{(n_1,v_1),(n_2,v_2),\dots,(n_N,v_N)},
    $$

    \noindent em que o par $(n_k,v_k)$ informa, simultaneamente, a ordem de visita do vértice $V_k$ e o veículo ao qual ele foi atribuído. Nessa representação, é conveniente introduzir o projetor:

    $$
    \Pi_k^{(a,\mu)}=\ket{a}\bra{a}_k^{(n)}\otimes \ket{\mu}\bra{\mu}_k^{(v)}
    $$

    \noindent que seleciona os estados nos quais o vértice $V_k$ ocupa a posição $a$ na rota do veículo $\mu$.

    Com essa nova codificação, o Hamiltoniano de custo pode ser generalizado de modo a contabilizar apenas deslocamentos entre vértices consecutivos atendidos pelo mesmo veículo. Em uma primeira formulação, esse termo pode ser escrito como:

    $$
        H_{\mathrm{cost}}^{(M)}
        =
        \sum_{\mu=1}^{M}
        \sum_{a=1}^{N-1}
        \sum_{i=1}^{N}
        \sum_{j=1}^{N}
        d_{ij}\,
        \left(\ket{a}\bra{a}_k^{(n)}\otimes \ket{\mu}\bra{\mu}_k^{(v)}\right)
        \left(\ket{a}\bra{a+1}_j^{(n)}\otimes \ket{\mu}\bra{\mu}_j^{(v)}\right)
    $$

    \noindent onde $d_{ij}$ representa o custo de deslocamento entre os vértices $V_i$ e $V_j$. Assim, o Hamiltoniano soma as distâncias entre clientes que ocupam posições consecutivas na rota de um mesmo veículo.

    Para assegurar a consistência combinatória da solução, é necessário impedir que dois clientes ocupem simultaneamente a mesma posição na rota de um mesmo veículo e, adicionalmente, evitar lacunas internas nas rotas construídas. Para isso, define-se o operador de ocupação
    $$
    \hat{N}_{a,\mu}=\sum_{k=1}^{N}\Pi_k^{(a,\mu)},
    $$

    \noindent a partir do qual se pode introduzir o termo de penalização
    $$
    H_{\mathrm{occ}}
    =
    \sum_{\mu=1}^{M}\sum_{a=1}^{N}
    \hat{N}_{a,\mu}\bigl(\hat{N}_{a,\mu}-1\bigr),
    $$

    \noindent que penaliza configurações em que mais de um cliente é associado à mesma posição $a$ do veículo $\mu$. Para desencorajar rotas com posições vazias intercaladas entre posições ocupadas, pode-se ainda acrescentar
    $$
    H_{\mathrm{gap}}
    =
    \sum_{\mu=1}^{M}\sum_{a=1}^{N-1}
    \bigl(1-\hat{N}_{a,\mu}\bigr)\hat{N}_{a+1,\mu},
    $$
    \noindent de modo que o Hamiltoniano total assuma a forma
    $$
    H_{\mathrm{VRP}}
    =
    H_{\mathrm{cost}}^{(M)}
    +\lambda_{1}H_{\mathrm{occ}}
    +\lambda_{2}H_{\mathrm{gap}},
    $$

    \noindent com $\lambda_1,\lambda_2>0$ parâmetros de penalização. Essa construção visa uma primeira generalização do modelo para o VRP de múltiplos veículos, preservando a estrutura em termos de projetores e abrindo caminho para a inclusão posterior de restrições adicionais, como capacidade e janelas de tempo.

    Como etapa posterior, pode-se incorporar ao formalismo uma restrição de capacidade por meio da introdução de uma variável de atribuição de veículo e de um operador de carga total associado a cada rota. Se $q_k$ denota a demanda do cliente $V_k$ e $Q_\mu$ a capacidade do veículo $\mu$, define-se inicialmente o operador
    $$
    \hat{y}_{k,\mu}=\sum_{a=1}^{N}\Pi_k^{(a,\mu)},
    $$
    o qual projeta sobre os estados em que o cliente $V_k$ é atendido pelo veículo $\mu$, independentemente de sua posição na rota. A carga total associada ao veículo $\mu$ pode então ser escrita como
    $$
    \hat{Q}_{\mu}=\sum_{k=1}^{N} q_k\,\hat{y}_{k,\mu}.
    $$

    Com isso, um possível Hamiltoniano de penalização para a restrição de capacidade é dado por

    $$
    H_{\mathrm{cap}}
    =
    \sum_{\mu=1}^{M}
    \Theta\!\left(\hat{Q}_{\mu}-Q_\mu\right)
    \left(\hat{Q}_{\mu}-Q_\mu\right)^2,
    $$

    \noindent em que $\Theta(\cdot)$ representa a função de Heaviside. Esse termo atribui energia positiva apenas às configurações em que a demanda total dos clientes atribuídos ao veículo $\mu$ excede sua capacidade máxima, preservando sem penalização os estados compatíveis com a restrição logística imposta.

    De modo análogo, a inclusão de janelas de tempo pode ser proposta mediante a introdução de operadores de tempo de chegada ao longo de cada rota. Se $[A_i,B_i]$ representa a janela de atendimento do cliente $V_i$, define-se $\hat{T}_{a,\mu}$ como o operador associado ao instante de chegada do veículo $\mu$ ao cliente visitado na posição $a$. A consistência temporal da rota pode ser imposta por meio do termo

    $$
    H_{\mathrm{temp}}
    =
    \sum_{\mu=1}^{M}\sum_{a=1}^{N-1}
    \left(
    \hat{T}_{a+1,\mu}
    -
    \hat{T}_{a,\mu}
    -
    \sum_{i=1}^{N}\sum_{j=1}^{N}
    (\tau_{ij}+s_i)\,
    \Pi_i^{(a,\mu)}\Pi_j^{(a+1,\mu)}
    \right)^2,
    $$

    \noindent em que $\tau_{ij}$ denota o tempo de deslocamento entre os clientes $V_i$ e $V_j$, e $s_i$ o tempo de serviço em $V_i$. A penalização associada às janelas de tempo pode então ser escrita como

    $$
    H_{\mathrm{jt}}
    =
    \sum_{\mu=1}^{M}\sum_{a=1}^{N}\sum_{i=1}^{N}
    \Pi_i^{(a,\mu)}
    \left[
    \Theta\!\left(A_i-\hat{T}_{a,\mu}\right)\left(A_i-\hat{T}_{a,\mu}\right)^2
    +
    \Theta\!\left(\hat{T}_{a,\mu}-B_i\right)\left(\hat{T}_{a,\mu}-B_i\right)^2
    \right].
    $$

    Nessa formulação, a energia do sistema cresce sempre que o atendimento ocorre antes da abertura ou após o fechamento da janela admissível, fornecendo, assim, uma proposta natural de extensão do Hamiltoniano para variantes do VRP com restrições temporais.

    Na etapa seguinte, os Hamiltonianos construídos serão empregados como funções de custo em modelos de otimização quântica variacional, com destaque para abordagens como o \textit{Variational Quantum Eigensolver} (VQE) e o \textit{Quantum Approximate Optimization Algorithm} (QAOA), adaptadas ao contexto dos Hamiltonianos formulados. Nessas abordagens, estados parametrizados serão preparados por circuitos variacionais compostos por portas gaussianas e, quando necessário, operações não gaussianas, sendo então avaliados iterativamente. Os parâmetros variacionais serão ajustados por rotinas de otimização clássica baseadas em gradiente ou métodos livres de derivadas, a depender da estrutura do Hamiltoniano e do comportamento numérico observado.

    A avaliação de desempenho não se restringirá aos modelos quânticos. Para fins de comparação, serão empregados métodos clássicos de referência, algoritmos exatos em instâncias de pequena dimensão. Essa abordagem permitirá analisar, de forma sistemática, a qualidade das soluções encontradas, o custo computacional envolvido, a estabilidade dos métodos frente à escolha de hiperparâmetros e o comportamento relativo entre abordagens clássicas e quânticas. Essa etapa constitui, portanto, não apenas uma validação prática do formalismo desenvolvido, mas também um estudo comparativo sobre o potencial e as limitações de diferentes estratégias de otimização aplicadas aos Hamiltonianos construídos.

    O cronograma de desenvolvimento do doutorado, posterior a etapa de qualificação, foi organizado de modo a articular, o aprofundamento do referencial teórico, o refinamento do modelo hamiltoniano proposto, sua implementação computacional a implementação de algoritmos de otimização quântica, validações necessárias e a redação progressiva da tese. Esse desenvolvimento está sintetizado na \autoref{tab:cronograma}. Busca-se, com isso, assegurar que a consolidação conceitual e a obtenção dos resultados numéricos avancem conjuntamente.

    \begin{table}[!htb]
    \centering
    \caption{Cronograma de execução das atividades}
    \label{tab:cronograma}
    \resizebox{\textwidth}{!}{
    \begin{tabular}{|p{6.2cm}|c|c|c|c|c|c|c|c|c|c|c|c|}
    \hline

    \textbf{Etapa} & \textbf{Jun} & \textbf{Jul} & \textbf{Ago} & \textbf{Set} & \textbf{Out} & \textbf{Nov} \\ \hline \hline
    Revisão bibliográfica e ampliação do referencial teórico &  \cellcolor{blue}  &  \cellcolor{blue}  &  \cellcolor{blue}  &     &     &     \\ \hline
    Desenvolvimento e extensão do Hamiltoniano               &  \cellcolor{green}  &  \cellcolor{green}  &  \cellcolor{green}  &  \cellcolor{green}  &     &     \\ \hline
    Implementação computacional  & &  \cellcolor{yellow}  &  \cellcolor{yellow}  &  \cellcolor{yellow}  &  \cellcolor{yellow}  &         \\ \hline
    Implementação computacional do algoritmo de otimização quântica & &  \cellcolor{blue}  &  \cellcolor{blue}  &  \cellcolor{blue}  &  \cellcolor{blue}  &         \\ \hline
    Simulações, análise de resultados e validação            &     &     &  \cellcolor{green}  &  \cellcolor{green}  &  \cellcolor{green}  &     \\ \hline
    Redação e consolidação da tese &  \cellcolor{yellow}  &  \cellcolor{yellow}  &  \cellcolor{yellow}  &  \cellcolor{yellow}  &  \cellcolor{yellow}  &  \cellcolor{yellow}  \\ \hline
    \end{tabular}
    }
    \end{table}